%------------------------------------------------------------------------------%
\documentclass[a4paper,12pt,fleqn]{article}

\usepackage{gaal}

\ShowAnswers

%------------------------------------------------------------------------------%
\begin{document}

\makeHeader{GAAL}{Prova 2 -- Turma 7}{23/09/2026}

% 01
%------------------------------------------------------------------------------%
\exercise{20\%}
Calcule o produto misto entre
\(
  \vec{p} = \vetortri{2}{1}{1}
\),
\(
  \vec{q} = \vetortri{1}{-1}{2}
\) e
\(
  \vec{r} = \vetortri{3}{2}{-1}
\)

\begin{answer}
  \begin{align*}
    \left(\vec{p} \times \vec{q}\right) \cdot \vec{r}
    & = \begin{vmatrix}
          2 & 1  & 1  \\
          1 & -1 & 2  \\
          3 & 2  & -1
        \end{vmatrix}
    \\
    & = 2 \begin{vmatrix}
        -1 & 2  \\
        2  & -1
      \end{vmatrix}
      - 1 \begin{vmatrix}
        1 & 2  \\
        3 & -1
      \end{vmatrix}
      + 1 \begin{vmatrix}
        1 & -1 \\
        3 & 2
      \end{vmatrix}
    \\
    & = 2(1-4) -(-1-6) + (2+3) \\
    & = -6 +7 +5 \\
    & = 6
  \end{align*}
\end{answer}

% 02
%------------------------------------------------------------------------------%
\exercise{20\%}
Determine o ângulo entre os vetores
\(
  \vec{v} = \vetortri{2}{1}{2}
\) e
\(
  \vec{w} = \vetortri{1}{-2}{2}
\)

\begin{answer}
  \[
    \vec{v} \cdot \vec{w}
    = \vetortri{2}{1}{2} \cdot \vetortri{1}{-2}{2}
    = 2 \times 1 + 1(-2) + 2 \times 2
    = 2 - 2 + 4
    = 4
  \]
  \[
    \norm{\vec{v}}
    = \norm{\vetortri{2}{1}{2}}
    = \sqrt{ 2^2 + 1^2 + 2^2}
    = \sqrt{9}
    = 3
  \]
  \[
    \norm{\vec{w}}
    = \norm{\vetortri{1}{-2}{2}}
    = \sqrt{ 1^2 + (-2)^2 + 2^2}
    = \sqrt{9}
    = 3
  \]
  \[
    \cos(\theta)
    = \frac{\vec{v}\cdot\vec{w}}{\norm{\vec{v}}\norm{\vec{w}}}
    = \frac{4}{9}
  \]
  \[
    \theta = \arccos\left(\frac{4}{9}\right)
    {\color{gray}
    \approx 1,11
    \approx \ang{63.61}
    }
  \]
\end{answer}

% 03
%------------------------------------------------------------------------------%
\exercise{20\%}
Determine a norma do vetor
\(
  \vec{\psi} = 3 \vetortri{1}{3}{-2} - 2 \vetortri{2}{-1}{-4}
\)

\begin{answer}
  \[
    \vec{\psi}
    = 3 \vetortri{1}{3}{-2} - 2 \vetortri{2}{-1}{-4}
    = \vetortri{3}{9}{-6} + \vetortri{-4}{2}{8}
    = \vetortri{-1}{11}{2}
  \]
  \[
    \norm{\vec{\psi}}
    = \norm{\vetortri{-1}{11}{2}}
    = \sqrt{(-1)^2+11^2+2^2}
    = \sqrt{1+121+4}
    = \sqrt{126}
  \]
\end{answer}

% 04
%------------------------------------------------------------------------------%
\exercise{20\%}
Calcule determinante de
\(
  A =
  \begin{pmatrix}
    2 & 1 & 1 & 0 \\
    1 & 2 & 0 & 1 \\
    1 & 0 & 2 & 1 \\
    0 & 1 & 1 & 2
  \end{pmatrix}
\)
e determine se a matriz é invertível

\begin{answer}
\begin{align*}
  \det(A)
  & = a_{11}(-1)^{1+1} M_{11}
    + a_{12}(-1)^{1+2} M_{12}
    + a_{13}(-1)^{1+3} M_{13}
    + a_{14}(-1)^{1+4} M_{14} \\
  & = 2 M_{11} - 1 M_{12} + 1 M_{13} - 0 M_{14}   \\
  & = 2 M_{11} - M_{12} + M_{13}
\end{align*}

\[
  M_{11}
  = \begin{vmatrix}
    2 & 0 & 1 \\
    0 & 2 & 1 \\
    1 & 1 & 2
  \end{vmatrix}
  = 2 (-1)^{1+1}
  \begin{vmatrix}
    2 & 1 \\
    1 & 2
  \end{vmatrix}
  + 0
  + 1 (-1)^{1+3}
  \begin{vmatrix}
    0 & 2 \\
    1 & 1
  \end{vmatrix}
  = 2(4-1)+(0-2)
  = 6 - 2
  = 4
\]
\[
  M_{12}
  = \begin{vmatrix}
    1 & 0 & 1 \\
    1 & 2 & 1 \\
    0 & 1 & 2
  \end{vmatrix}
  = 1 (-1)^{1+1}
  \begin{vmatrix}
    2 & 1 \\
    1 & 2
  \end{vmatrix}
  + 0
  + 1 (-1)^{1+3}
  \begin{vmatrix}
    1 & 2 \\
    0 & 1
  \end{vmatrix}
  = (4-1)+(1-0)
  = 3 + 1
  = 4
\]
\[
  M_{13}
  = \begin{vmatrix}
    1 & 2 & 1 \\
    1 & 0 & 1 \\
    0 & 1 & 2
  \end{vmatrix}
  = 1 (-1)^{2+1}
  \begin{vmatrix}
    1 & 1 \\
    1 & 2
  \end{vmatrix}
  + 0
  + 1 (-1)^{2+3}
  \begin{vmatrix}
    1 & 2 \\
    0 & 1
  \end{vmatrix}
  = -(2-1) - (1-0)
  = -1
\]
\[
  \det(A)
  = 2 \times 4 - 4 + (-1)
  = 3
\]

Como
\(
  \det(A) \neq 0
\)
$A$ é invertível
\end{answer}

% 05
%------------------------------------------------------------------------------%
\exercise{20\%}
Encontre o vetor unitário na direção e sentido de $\vec{p}\times\vec{q}$ com
$
  \vec{p} = \vetortri{2}{1}{-1}
$ e
$
  \vec{q} = \vetortri{1}{-2}{3}
$
\begin{answer}
  \begin{align*}
    \vec{\psi}
    & = \vec{p}\times\vec{q} \\
    & = \begin{vmatrix}
        \vec{i} & \vec{j} & \vec{k} \\
        2       & 1       & -1      \\
        1       & -2      & 3
      \end{vmatrix}
    \\
    & = \vec{i} (-1)^{1+1}
      \begin{vmatrix}
        1  & -1 \\
        -2 & 3
      \end{vmatrix}
      + \vec{j} (-1)^{1+2}
      \begin{vmatrix}
        2 & -1 \\
        1 & 3
      \end{vmatrix}
      + \vec{k} (-1)^{2+2}
      \begin{vmatrix}
        2 & 1 \\
        1 & -2
      \end{vmatrix}
    \\
    & = \vec{i} ( 3 - 2)
      - \vec{j} ( 6 + 1)
      + \vec{k} (-4 - 1)
    \\
    & = \vec{i} - 7\vec{j} -5 \vec{k} \\
    & = \vetortri{1}{-7}{-5}
  \end{align*}
  \[
    \norm{\vec{\psi}}
    = \norm{\vetortri{1}{-7}{-5}}
    = \sqrt{1^2 + (-7)^2 + (-5)^2}
    = \sqrt{1 + 49 + 25}
    = \sqrt{75}
  \]
  \[
    \hat{\psi}
    = \frac{\vec{\psi}}{\norm{\vec{\psi}}}
    = \frac{1}{\sqrt{75}} \vetortri{1}{-7}{-5}
    = \Vetortri{\frac{1}{\sqrt{75}}}{\frac{-7}{\sqrt{75}}}{\frac{-5}{\sqrt{75}}}
  \]

\end{answer}

%------------------------------------------------------------------------------%
\end{document}
%------------------------------------------------------------------------------%

